%! TEX program = xelatex
%! TEX root = ../main.tex
%! TEX encoding = utf-8

\chapter{水下快速测距算法研究与仿真}[Fast Underwater Ranging Algorithm and Simulation]
\label{chap04}

本章只讨论测距问题，形成“往返时延模型—时间戳提取—声速与声线模型—快速距离反演—误差分析”的独立证据链。通信算法仅作为请求和应答的承载方式，不在本章重复讨论SFH-MFSK调制与解调。

\section{往返时延测量模型}[Round-Trip Delay Measurement Model]

设甲板单元发射请求的时刻为$t_1$，水下机检测到请求首达径的时刻为$t_2$，水下机完成处理并发射应答的时刻为$t_3$，甲板单元检测到应答首达径的时刻为$t_4$。则
\begin{equation}
 \Delta t_{\rm meas}=t_4-t_1
 =\tau_{\rm up}+\tau_{\rm proc}+\tau_{\rm down},
\label{eq:chap04_rtt}
\end{equation}
其中$\tau_{\rm proc}=t_3-t_2$为水下机固定处理延时。经标定得到$\hat{\tau}_{\rm proc}$后，若上下行传播路径可视为互易，则单程传播时延估计为
\begin{equation}
 \hat{\tau}_{\rm prop}
 =\frac{\Delta t_{\rm meas}-\hat{\tau}_{\rm proc}}{2}.
\label{eq:chap04_prop}
\end{equation}
式\eqref{eq:chap04_prop}是声线反演的时间输入，不等同于距离。处理延时标定误差、时间戳量化误差和上下行传播不对称均会进入最终测距误差。

\section{首达径检测与时间戳提取}[First-Arrival Detection and Timestamp Extraction]

多径条件下，最大相关峰可能属于后续反射路径，因此采用匹配滤波和首达径门限联合判决。设发送参考信号为$s(t)$，接收信号为$r(t)$，匹配滤波输出为
\begin{equation}
 y(\tau)=\int r(t)s(t-\tau)\mathrm{d}t .
\label{eq:chap04_match}
\end{equation}
离散实现中，在噪声估计基础上设置门限$\gamma$。候选峰$y[n]$满足
\begin{equation}
 |y[n]|>\gamma
\label{eq:chap04_threshold}
\end{equation}
且其前方保护区间不存在更早的高可信峰时，记录其位置为首达径时间戳。门限、保护区间和最小峰间隔必须在仿真与平台测试中保持一致。

\section{声速剖面与声线传播模型}[Sound-Speed Profile and Ray Model]

声速由温度、盐度和深度共同决定，可采用
\begin{equation}
 c(T,S,z)=1449.2+4.6T-0.055T^2+2.9\times10^{-4}T^3
 +(1.34-0.01T)(S-35)+0.016z
\label{eq:chap04_sound_speed}
\end{equation}
进行工程估计。对于水平分层介质，设声线与水平面的夹角为$\theta(z)$，Snell定律写为
\begin{equation}
 \frac{\cos\theta(z)}{c(z)}=p,
\label{eq:chap04_snell}
\end{equation}
其中$p$为水平慢度。无转向点时，候选慢度对应的水平距离和传播时间为
\begin{equation}
\begin{aligned}
 X(p)&=\int_{z_1}^{z_2}
 \frac{pc(z)}{\sqrt{1-p^2c^2(z)}}\,\mathrm dz,\\
 T(p)&=\int_{z_1}^{z_2}
 \frac{1}{c(z)\sqrt{1-p^2c^2(z)}}\,\mathrm dz.
\end{aligned}
\label{eq:chap04_ray}
\end{equation}
存在转向点时，应按传播段分别计算并累加。

\section{不依赖目标深度的快速距离反演}[Depth-Independent Fast Distance Inversion]

仅凭单次传播时延和声速剖面，通常不能同时唯一确定距离和未知端点深度。因此，本文所称“不依赖目标深度”不是删除所有几何信息，而是不将目标深度作为独立人工输入；算法必须使用项目中已知的几何约束或其他可观测量。定稿时应明确列出约束量、未知量和可解条件。

在约束关系已确定后，构造时间残差
\begin{equation}
 e(p)=T(p)-\hat{\tau}_{\rm prop}.
\label{eq:chap04_residual}
\end{equation}
在有效慢度区间$[p_{\rm L},p_{\rm U}]$内，若残差连续且端点异号，可采用二分搜索：
\begin{equation}
 p_j=\frac{p_{{\rm L},j}+p_{{\rm U},j}}{2}.
\label{eq:chap04_bisection}
\end{equation}
根据$e(p_j)$的符号更新区间，直至$|e(p_j)|<\varepsilon_t$或达到最大迭代次数。收敛后由$X(p_j)$输出距离。若存在多个声线分支，应根据首达径和已知几何约束选取解，并输出残差、迭代次数和失败标志。若约束不足以保证唯一解，则应将该工况标记为不可解，而不是输出成功形态的距离。

\section{测距仿真条件与评价指标}[Ranging Simulation Conditions and Metrics]

测距仿真应分别设置不同真实距离、信噪比、多径强度、声速剖面和首达径检测误差。主要指标包括平均误差、偏差、均方根误差、最大绝对误差、不可解比例和单次计算时间：
\begin{equation}
 {\rm RMSE}=\sqrt{\frac{1}{N}\sum_{i=1}^{N}(\hat R_i-R_i)^2}.
\label{eq:chap04_rmse}
\end{equation}
对比方案至少应包含常声速模型、等效声速修正和快速声线修正，并报告每种方案使用的输入信息，确保“不依赖目标深度”的比较成立。

\begin{table}[htbp]
 \bicaption[tab:chap04_params]{}{测距仿真参数}{Tab.}{Ranging simulation parameters}
 \vspace{0.5em}\centering\wuhao
 \begin{tabular}{ccc}
 \toprule[1.5pt]
 参数 & 符号 & 取值\\
 \midrule[1pt]
 真实距离范围 & $R$ & 待填写\\
 信噪比范围 & $\mathrm{SNR}$ & 待填写\\
 声速剖面 & $c(z)$ & 待填写\\
 首达径误差 & $\delta t$ & 待填写\\
 慢度搜索容差 & $\varepsilon_t$ & 待填写\\
 最大迭代次数 & $J_{\max}$ & 待填写\\
 \bottomrule[1.5pt]
 \end{tabular}
\end{table}

\section{仿真结果与误差分析}[Simulation Results and Error Analysis]

当前材料未提供完整测距仿真输出，故本节保留结果填写边界。定稿时应先报告处理延时已补偿条件下的基准结果，再分别改变首达径误差、声速剖面和多径条件，最后给出计算效率。误差分析应区分首达径检测误差、处理延时标定误差、声速模型误差、声线分支选择错误和几何约束误差；不能将所有误差归因于声速修正。

\section{本章小结}[Summary]

本章建立了基于请求—应答往返时延的测距模型，给出了处理延时补偿、首达径检测、声速剖面建模和声线距离反演方法，并明确了不依赖目标深度算法必须满足的几何约束与可解条件。后续仿真将以误差统计和计算效率为核心，所得基准结果作为第5章平台测距验证的对照。
